%%%%%%%%%%%%%%%%
%%% gen 50 %%%
%%%%%%%%%%%%%%%%

\Chapter{50}

\Verse{1}
{
	\hJ{וַיִּפֹּל יוֹסֵף עַל פְּנֵי אָבִיו וַיֵּבְךְּ עָלָיו וַיִּשַּׁק לוֹ׃}
}
{
	\eJ{
		50 And Joseph fell upon his father’s face and wept over
		him and kissed him.
	}
}
{
}

\Verse{2}
{
	\hJ{וַיְצַו יוֹסֵף אֶת עֲבָדָיו אֶת הָרֹפְאִים לַחֲנֹט אֶת אָבִיו וַיַּחַנְטוּ הָרֹפְאִים אֶת יִשְׂרָאֵל׃}
}
{
	\eJ{
		And Joseph commanded his servants, the physicians, to
		embalm his father. And the physicians embalmed Israel,
	}
}
{
}

\Verse{3}
{
	\hJ{וַיִּמְלְאוּ לוֹ אַרְבָּעִים יוֹם כִּי כֵּן יִמְלְאוּ יְמֵי הַחֲנֻטִים וַיִּבְכּוּ אֹתוֹ מִצְרַיִם שִׁבְעִים יוֹם׃}
}
{
	\eJ{
		and they spent forty days on him, because that is the
		number of days they spent on the embalming, and Egypt
		mourned him seventy days.
	}
}
{
}

\Verse{4}
{
	\hJ{וַיַּעַבְרוּ יְמֵי בְכִיתוֹ וַיְדַבֵּר יוֹסֵף אֶל בֵּית פַּרְעֹה לֵאמֹר אִם נָא מָצָאתִי חֵן בְּעֵינֵיכֶם דַּבְּרוּ נָא בְּאזְנֵי פַרְעֹה לֵאמֹר׃}
}
{
	\eJ{
		And his mourning days passed, and Joseph spoke to
		Pharaoh’s house, saying, “If I’ve found favor in your
		eyes, speak in Pharaoh’s ears, saying,
	}
}
{
}

\Verse{5}
{
	\hJ{אָבִי הִשְׁבִּיעַנִי לֵאמֹר הִנֵּה אָנֹכִי מֵת בְּקִבְרִי אֲשֶׁר כָּרִיתִי לִי בְּאֶרֶץ כְּנַעַן שָׁמָּה תִּקְבְּרֵנִי וְעַתָּה אֶעֱלֶה נָּא וְאֶקְבְּרָה אֶת אָבִי וְאָשׁוּבָה׃}
}
{
	\eJ{
		‘My father had me swear, saying, “Here I’m dying. In my
		tomb, which I dug for me in the land of Canaan: you shall
		bury me there.” And now, let me go up so I may bury my
		father and come back.’”
	}
}
{
}

\Verse{6}
{
	\hJ{וַיֹּאמֶר פַּרְעֹה עֲלֵה וּקְבֹר אֶת אָבִיךָ כַּאֲשֶׁר הִשְׁבִּיעֶךָ׃}
}
{
	\eJ{
		And Pharaoh said, “Go up and bury your father as he had
		you swear.”
	}
}
{
}

\Verse{7}
{
	\hJ{וַיַּעַל יוֹסֵף לִקְבֹּר אֶת אָבִיו וַיַּעֲלוּ אִתּוֹ כּל עַבְדֵי פַרְעֹה זִקְנֵי בֵיתוֹ וְכֹל זִקְנֵי אֶרֶץ מִצְרָיִם׃}
}
{
	\eJ{
		And Joseph went up to bury his father, and all of
		Pharaoh’s servants, the elders of his house, and all the
		elders of the land of Egypt,
	}
}
{
}

\Verse{8}
{
	\hJ{וְכֹל בֵּית יוֹסֵף וְאֶחָיו וּבֵית אָבִיו רַק טַפָּם וְצֹאנָם וּבְקָרָם עָזְבוּ בְּאֶרֶץ גֹּשֶׁן׃}
}
{
	\eJ{
		and all of Joseph’s house and his brothers and his
		father’s house. Only their infants and their flock and
		their oxen they left in the land of Goshen.
	}
}
{
}

\Verse{9}
{
	\hJ{וַיַּעַל עִמּוֹ גַּם רֶכֶב גַּם פָּרָשִׁים וַיְהִי הַמַּחֲנֶה כָּבֵד מְאֹד׃}
}
{
	\eJ{
		And both chariots and horsemen went up with him. It was a
		very heavy camp.
	}
}
{
}

\Verse{10}
{
	\hJ{וַיָּבֹאוּ עַד גֹּרֶן הָאָטָד אֲשֶׁר בְּעֵבֶר הַיַּרְדֵּן וַיִּסְפְּדוּ שָׁם מִסְפֵּד גָּדוֹל וְכָבֵד מְאֹד וַיַּעַשׂ לְאָבִיו אֵבֶל שִׁבְעַת יָמִים׃}
}
{
	\eJ{
		And they came to the threshing floor of Atad, which is
		across the Jordan, and they had a very big and heavy
		funeral, and he made a mourning for his father or seven
		days.
	}
}
{
}

\Verse{11}
{
	\hJ{וַיַּרְא יוֹשֵׁב הָאָרֶץ הַכְּנַעֲנִי אֶת הָאֵבֶל בְּגֹרֶן הָאָטָד וַיֹּאמְרוּ אֵבֶל כָּבֵד זֶה לְמִצְרָיִם עַל כֵּן קָרָא שְׁמָהּ אָבֵל מִצְרַיִם אֲשֶׁר בְּעֵבֶר הַיַּרְדֵּן׃}
}
{
	\eJ{
		And the Canaanite residents of the land saw the mourning
		at the threshing floor of Atad and said, “This is a heavy
		mourning to Egypt.” On account of this its name was called
		Abel of Egypt, which is beyond the Jordan.
	}
}
{
}

\Verse{12}
{
	\hP{וַיַּעֲשׂוּ בָנָיו לוֹ כֵּן כַּאֲשֶׁר צִוָּם׃}
}
{
	\eP{And his sons did so for him, as he had commanded them.}
}
{
}

\Verse{13}
{
	\hP{וַיִּשְׂאוּ אֹתוֹ בָנָיו אַרְצָה כְּנַעַן וַיִּקְבְּרוּ אֹתוֹ בִּמְעָרַת שְׂדֵה הַמַּכְפֵּלָה אֲשֶׁר קָנָה אַבְרָהָם אֶת הַשָּׂדֶה לַאֲחֻזַּת קֶבֶר מֵאֵת עֶפְרֹן הַחִתִּי עַל פְּנֵי מַמְרֵא׃}
}
{
	\eP{
		And his sons carried him to the land of Canaan and buried
		him in the cave of the field of Machpelah, as Abraham had
		bought the field as a possession for a tomb from Ephron,
		the Hittite, facing Mamre.
	}
}
{
}

\Verse{14}
{
	\hJ{וַיָּשׁב יוֹסֵף מִצְרַיְמָה הוּא וְאֶחָיו וְכל הָעֹלִים אִתּוֹ לִקְבֹּר אֶת אָבִיו אַחֲרֵי קבְרוֹ אֶת אָבִיו׃}
}
{
	\eJ{
		And, after his burial of his father, Joseph came back to
		Egypt, he and his brothers and all those who went up with
		him to bury his father.
	}
}
{
}

\Verse{15}
{
	\hE{וַיִּרְאוּ אֲחֵי יוֹסֵף כִּי מֵת אֲבִיהֶם וַיֹּאמְרוּ לוּ יִשְׂטְמֵנוּ יוֹסֵף וְהָשֵׁב יָשִׁיב לָנוּ אֵת כּל הָרָעָה אֲשֶׁר גָּמַלְנוּ אֹתוֹ׃}
}
{
	\eE{
		And Joseph’s brothers saw that their father was dead, and
		they said, “If Joseph will despise us he’ll pay us back
		all the bad that we dealt to him.”
	}
}
{
}

\Verse{16}
{
	\hE{וַיְצַוּוּ אֶל יוֹסֵף לֵאמֹר אָבִיךָ צִוָּה לִפְנֵי מוֹתוֹ לֵאמֹר׃}
}
{
	\eE{
		And they commanded to Joseph, saying, “Your father had
		commanded before his death, saying,
	}
}
{
%	\footnote{
%		they commanded. Why would they command? The problem appears to be
%		scribal. The Septuagint appears to be translating Hebrew , meaning “they
%		approached,” instead of the MT , meaning “they commanded.” A scribe
%		apparently changed the word under the influence of the word “command”
%		which comes up later in the verse. Footnote 50:16. Your father had
%		commanded before his death. We never find out whether Joseph—or his
%		brothers—ever told Jacob what his brothers did to him. The brothers
%		claim that Jacob commanded that Joseph should forgive them, but we do
%		not know if they are making this up or not. Either way, it is the right
%		message: after a parent’s death, the children should try to heal any old
%		wounds and draw close.
%	}
}

\Verse{17}
{
	\hE{כֹּה תֹאמְרוּ לְיוֹסֵף אָנָּא שָׂא נָא פֶּשַׁע אַחֶיךָ וְחַטָּאתָם כִּי רָעָה גְמָלוּךָ וְעַתָּה שָׂא נָא לְפֶשַׁע עַבְדֵי אֱלֹהֵי אָבִיךָ וַיֵּבְךְּ יוֹסֵף בְּדַבְּרָם אֵלָיו׃}
}
{
	\eE{
		‘You shall say this to Joseph: Please, bear your brothers’
		offense and their sin, because they dealt you bad.’ And
		now bear the offense of the servants of your father’s
		{\God}.” And Joseph wept when they spoke to him.
	}
}
{
}

\Verse{18}
{
	\hE{וַיֵּלְכוּ גַּם אֶחָיו וַיִּפְּלוּ לְפָנָיו וַיֹּאמְרוּ הִנֶּנּוּ לְךָ לַעֲבָדִים׃}
}
{
	\eE{
		And his brothers also went and fell in front of him and
		said, “Here, we’re yours as slaves.”
	}
}
{
%	\footnote{
%		we’re yours as slaves. This is the final ironic recompense in the chain
%		of deceptions in this family. The brothers who once sold Joseph as a
%		slave now say that they will be his slaves.
%	}
}

\Verse{19}
{
	\hE{וַיֹּאמֶר אֲלֵהֶם יוֹסֵף אַל תִּירָאוּ כִּי הֲתַחַת אֱלֹהִים אָנִי׃}
}
{
	\eE{
		And Joseph said to them, “Don’t be afraid, because am I in
		{\God}’s place?
	}
}
{
}

\Verse{20}
{
	\hE{וְאַתֶּם חֲשַׁבְתֶּם עָלַי רָעָה אֱלֹהִים חֲשָׁבָהּ לְטֹבָה לְמַעַן עֲשֹׂה כַּיּוֹם הַזֶּה לְהַחֲיֹת עַם רָב׃}
}
{
	\eE{
		And you thought bad against me. {\God} thought for good: in
		order to do as it is today, to keep alive a numerous
		people.
	}
}
{
}

\Verse{21}
{
	\hE{וְעַתָּה אַל תִּירָאוּ אָנֹכִי אֲכַלְכֵּל אֶתְכֶם וְאֶת טַפְּכֶם וַיְנַחֵם אוֹתָם וַיְדַבֵּר עַל לִבָּם׃}
}
{
	\eE{
		And now, don’t be afraid. I’ll provide for you and your
		infants.” And he consoled them. And he spoke on their
		heart.
	}
}
{
%	\footnote{
%		he spoke on their heart. We have seen that each act of deception since
%		Jacob led to another deception that came as a recompense. Thus
%		deceptions and hurts within a family can go on in a perpetual cycle. In
%		order to bring it to an end, one member of the family who is entitled to
%		retribution must stop the cycle and forgive instead. That is what Joseph
%		does here. This beautiful end of the story of deceptions and
%		retributions completes an intricately embroidered narrative. Many
%		interpreters have treated the Joseph cycle as if it were a separate
%		entity, referring to it as a “novella.” Insofar as that term connotes an
%		independent work, it is misleading, cutting off what is an integral part
%		of this connected narrative. Symbols, ironies, and relationships of
%		characters and events are lost when one performs this surgery and
%		separates the Joseph cycle from the rest of the story.
%	}
}

\Verse{22}
{
	\hJ{וַיֵּשֶׁב יוֹסֵף בְּמִצְרַיִם הוּא וּבֵית אָבִיו וַיְחִי יוֹסֵף מֵאָה וָעֶשֶׂר שָׁנִים׃}
}
{
	\eJ{
		And Joseph lived in Egypt, he and his father’s house. And
		Joseph lived a hundred ten years,
	}
}
{
}

\Verse{23}
{
	\hE{וַיַּרְא יוֹסֵף לְאֶפְרַיִם בְּנֵי שִׁלֵּשִׁים גַּם בְּנֵי מָכִיר בֶּן מְנַשֶּׁה יֻלְּדוּ עַל בִּרְכֵּי יוֹסֵף׃}
}
{
	\eE{
		and Joseph saw children of Ephraim’s of the third
		generation. Also children of Machir, son of Manasseh, were
		born on Joseph’s knees.
	}
}
{
}

\Verse{24}
{
	\hE{וַיֹּאמֶר יוֹסֵף אֶל אֶחָיו אָנֹכִי מֵת וֵאלֹהִים פָּקֹד יִפְקֹד אֶתְכֶם וְהֶעֱלָה אֶתְכֶם מִן הָאָרֶץ הַזֹּאת אֶל הָאָרֶץ אֲשֶׁר נִשְׁבַּע לְאַבְרָהָם לְיִצְחָק וּלְיַעֲקֹב׃}
}
{
	\eE{
		And Joseph said to his brothers, “I’m dying. And {\God} will
		take account of you and will bring you up from this land
		to the land that He swore to Abraham, to Isaac, and to
		Jacob.”
	}
}
{
}

\Verse{25}
{
	\hE{וַיַּשְׁבַּע יוֹסֵף אֶת בְּנֵי יִשְׂרָאֵל לֵאמֹר פָּקֹד יִפְקֹד אֱלֹהִים אֶתְכֶם וְהַעֲלִתֶם אֶת עַצְמֹתַי מִזֶּה׃}
}
{
	\eE{
		And Joseph had the children of Israel swear, saying, “{\God}
		will take account of you, and you’ll bring up my bones
		from here.”
	}
}
{
}

\Verse{26}
{
	\hE{וַיָּמת יוֹסֵף בֶּן מֵאָה וָעֶשֶׂר שָׁנִים וַיַּחַנְטוּ אֹתוֹ וַיִּישֶׂם בָּאָרוֹן בְּמִצְרָיִם׃}
}
{
	\eE{
		And Joseph died, a hundred ten years old, and they
		embalmed him and set him in a coffin in Egypt. COMMENTARY
		ON THE TORAH The book of Exodus tells the story of the
		birth of a nation in slavery and ends with the nation’s
		establishment of its own center, leaders, and symbols in
		freedom. Genesis involves a continuous narrowing of
		attention from the universe to the earth to humanity to a
		particular family; Exodus begins to broaden the
		circumference of attention again as the family grows into
		a nation—and comes into conflict with another nation.
		Whereas Genesis sets the rest of the books of the Bible in
		context, Exodus does not set them in context so much as
		introduce fundamental components that will function
		centrally in almost all the coming books of the Bible.
		Exodus introduces the nation of Israel. It introduces
		prophecy. It introduces law. Arguably most important of
		all, it introduces the theme of YHWH’s becoming known to
		the world.
	}
}
{
}
